\chapter{Labor Signs (Signs of Labor)}

\section*{Definition}
Labor signs are the indicators that a pregnant woman is entering the active phase of childbirth. Recognizing these signs helps determine when to seek medical care.

\section*{What Happens in Your Body}
As pregnancy reaches term (37-42 weeks), hormonal changes trigger uterine contractions, cervical ripening, and effacement. Prostaglandins and oxytocin stimulate coordinated uterine contractions. The cervix softens, shortens (effaces), and opens (dilates). The presenting fetal part descends into the pelvis. The amniotic sac may rupture spontaneously.

\section*{Causes}
\begin{itemize}
\item Natural progression of full-term pregnancy
\item Preterm labor occurs before 37 weeks
\item Induced labor (medical or elective)
\item Spontaneous rupture of membranes
\end{itemize}

\section*{When to See a Doctor}
\begin{itemize}
\item Regular, painful contractions that become stronger and closer together
\item Contractions lasting 30-70 seconds and occurring every 5-10 minutes
\item Lower back pain that does not go away
\item Bloody show: pink or blood-tinged mucus discharge
\item Rupture of membranes: gush or leak of fluid from the vagina
\item Water breaking may feel like a slow trickle or sudden gush
\item Cervical dilation and effacement on examination
\end{itemize}

\section*{Related Symptoms}
\begin{itemize}
\item False labor (Braxton Hicks contractions, which are irregular and do not dilate the cervix)
\item Preterm labor (before 37 weeks)
\item Vaginal bleeding in pregnancy
\item Premature rupture of membranes
\item Back pain during pregnancy
\end{itemize}

\section*{Self-Care / Home Management}
\begin{itemize}
\item Time contractions from the start of one to the start of the next
\item Drink fluids and rest between contractions
\item Use breathing and relaxation techniques
\item Take a warm shower or bath for comfort
\item Change positions frequently (walking, sitting, lying on side)
\item Eat light snacks if hungry (avoid heavy meals)
\item Call the healthcare provider or go to the hospital when contractions are 5 minutes apart lasting 1 minute for 1 hour
\end{itemize}

\section*{How Your Doctor Will Evaluate You}
\begin{itemize}
\item Pelvic examination to check for cervical dilation and effacement
\item Fetal monitoring for contraction pattern and fetal heart rate
\item AmniSure or Nitrazine test to confirm rupture of membranes
\item Ultrasound to assess fetal position, fluid volume, and estimated weight
\end{itemize}

\section*{Treatment Options}
\begin{itemize}
\item Admit for active labor management
\item Pain management: epidural, IV narcotics, nitrous oxide
\item Continuous fetal monitoring
\item Oxytocin augmentation if labor slows
\item Assisted delivery (forceps, vacuum) if needed
\item Cesarean section for complications
\item Neonatal care support after delivery
\end{itemize}

\section*{References}
\begin{thebibliography}{99}
\bibitem{labor_signs:premature1} American College of Obstetricians and Gynecologists. "Management of preterm labor." ACOG Practice Bulletin No. 171. Obstet Gynecol. 2016;128(4):e155-e164.
\bibitem{labor_signs:premature2} Romero R, Dey SK, Fisher SJ. "Preterm labor: one syndrome, many causes." Science. 2014;345(6198):760-765.
\bibitem{labor_signs:premature3} Goldenberg RL, Culhane JF, Iams JD, et al. "Epidemiology and causes of preterm birth." Lancet. 2008;371(9606):75-84.
\bibitem{labor_signs:premature4} Iams JD. "Preterm labor and birth." In: Williams Obstetrics. 26th ed. McGraw-Hill; 2022.
\bibitem{labor_signs:premature5} Haas DM, Caldwell DM, Kirkpatrick P, et al. "Tocolytic therapy for preterm delivery: systematic review and network meta-analysis." BMJ. 2012;345:e6226.
\bibitem{labor_signs:premature6} Mercer BM, Goldenberg RL, Das A, et al. "The preterm prediction study: a clinical risk assessment system." Am J Obstet Gynecol. 1996;174(6):1885-1893.
\end{thebibliography}
